% !TeX root = ../complete.tex
\chapter*{جمع‌بندی جزوه}
\phantomsection
\addcontentsline{toc}{chapter}{جمع‌بندی جزوه}

تحلیل مدار \lr{CMOS} از اتصال‌های مدار آغاز می‌شود. مرجع ولتاژ، جهت
جریان و وضعیت هر ورودی را مشخص کنید. با \hyperref[sec:book-cmos]{مدل
\lr{Switch}}، مسیر اتصال خروجی به تغذیه یا زمین و مقدار منطقی مورد
انتظار را به دست آورید. همان مسیر را در \hyperref[sec:book-fabrication]{\lr{Layout}}
دنبال کنید تا ارتباط طرح فیزیکی با توپولوژی مدار روشن شود.

برای تحلیل کمّی، ولتاژ پایانه‌ها را تعیین و ناحیه کاری هر افزاره را
فرض کنید. از \hyperref[sec:book-mos-current]{مدل جریان \lr{Long-channel}}
آغاز کنید و متناسب با بایاس و هندسه، \hyperref[sec:nonideal-current]{اثرهای
غیرایده‌آل} را وارد کنید. در مسئله زمانی، \hyperref[sec:book-cv]{ظرفیت‌ها}
و تغییر ولتاژ پایانه‌ها نیز مقدار بار جابه‌جاشده را تعیین می‌کنند.

برای یافتن خروجی ماندگار، تعادل جریان گره را بنویسید. پس از حل، ناحیه کاری،
علامت جریان، واحد و دامنه ولتاژ پاسخ را بررسی کنید. با تکرار این تحلیل برای
ورودی‌های مختلف، \hyperref[sec:dc-transfer]{مشخصه انتقال \lr{DC}} حاصل
می‌شود. از روی این مشخصه می‌توان \hyperref[sec:noise-margin]{سطوح معتبر
منطقی و \lr{Noise Margin}} را به دست آورد. در \hyperref[sec:pass-transistor]{مدار
\lr{Pass Transistor}}، همین کنترل ناحیه کاری روشن می‌کند که خروجی تا کدام
سطح ولتاژ می‌تواند تغییر کند.
